\begin{center}\bfseries 8.\end{center}
\noindent\hspace*{\parindent}\textsc{Gauss} errichtet nun das Gebäude seiner Theorie auf zwei algebraischen Grundeigenschaften der quadratischen Formen.

Die eine von ihnen ist die Identität
\[
f(x,y)\mathbin{.}f(x',y')=[(ax+by)x'+(bx+cy)y']^2-D(xy'-x'y)^2,
\]
worin \(D\), die Verknüpfung der Koeffizienten
\[
D=b^2-ac,
\]
als eine die ganze Theorie beherrschende Zahl von \textsc{Gauss} als Determinante der Form benannt worden ist; je nach ihrem Vorzeichen zerfallen die Formen in zwei Gattungen, welche vielfach ein sehr verschiedenes Verhalten zeigen.
\newpage\SourceLeaf{26}\ArithmeticHead{19}{R}{BINÄRE QUADRATISCHE FORMEN.}\HistoricalFoot{}{3*}
\noindent Aus der obigen Identität ergibt sich sogleich die zur eigentlichen Darstellung von \(n\) durch \((a,b,c)\) d. i. zur Darstellung mittels teilerfremder Werte \(x,\,y\) notwendige Bedingung, dass \(D\) quadratischer Rest von \(n\) sei, sowie die Erkenntnis, dass zwischen den vorhandenen eigentlichen Darstellungen und den Wurzeln der Kongruenz
\[
z^2\equiv D\ (\mathrm{mod.}\ n)
\]
ein enger Zusammenhang besteht, dem gemäss jede solche Darstellung als zu einer bestimmten der Kongruenzwurzeln »gehörig« bezeichnet werden darf.

Die zweite Grundeigenschaft ist die Transformation einer quadratischen Form
\[
ax^2+2bxy+cy^2
\]
mittels einer linearen ganzzahligen Substitution
\[
x=\alpha x'+\beta y',\ y=\gamma x'+\delta y'
\]
in eine andere Form
\[
a'x'^2+2b'x'y'+c'y'^2,
\]
deren Determinante
\[
D'=D\mathbin{.}(\alpha\delta-\beta\gamma)^2
\]
ist und welche unter der ersteren enthalten, insbesondere, wenn der Substitutionsmodul \(\alpha\delta-\beta\gamma=\pm1\) ist, ihr äquivalent genannt wird. Die zuerst von \textsc{Gauss} getroffene Unterscheidung solcher Äquivalenz in eine eigentliche und uneigentliche je nach dem Vorzeichen von \(\pm1\) verleiht der ferneren Entwicklung erhebliche Klarheit und Einfachheit.

Das Problem der Darstellung einer Zahl \(n\) durch eine gegebene Form \(f\) wird nunmehr zurückgeführt auf die Frage, ob diese Form mit gewissen andern Formen \(F\) mit gleicher Determinante, von denen je eine der Wurzel entspricht, zu der die Darstellung gehörig gedacht wird, äquivalent sei. Die neue Frage aber entscheidet sich durch eine zwiefache Untersuchung. Die eine von ihnen lehrt, dass eine gegebene Form durch eine Kette von sogenannten benachbarten Formen, deren jede mit der vorhergehenden eigentlich äquivalent ist, in eine ihr eigentlich äquivalente Form übergeführt werden kann, deren Koeffizienten durch gewisse Ungleichheiten beschränkt sind und welche reduziert genannt wird,
\newpage\SourceLeaf{27}\ArithmeticHead{20}{L}{BACHMANN, ÜBER GAUSS’ ZAHLENTHEORETISCHE ARBEITEN.}
\noindent womit sich dann zugleich auch eine Transformation der gegebenen Form in die reduzierte ergibt. Diesen Umstand und dass die Anzahl der reduzierten Formen nur endlich ist, fand \textsc{Gauss} schon als von \textsc{Lagrange}\footnote{Nouv. Mém. de l’Acad. Berlin (1773) 1775, S. 265; Oeuvres III, S. 695.}\ nachgewiesen vor. Durch die andere Untersuchung wird festgestellt, ob zwei solche reduzierte Formen selbst einander äquivalent sind oder nicht, und im erstern Falle eine Transformation der einen in die andere gegeben. Hierdurch ermöglicht sich nicht nur die Entscheidung jener Frage für jede der gedachten Formen \(F\), sondern auch im Falle der Äquivalenz eine Transformation der gegebenen Form in jede ihr äquivalente unter diesen.

Die Bestimmung der reduzierten Formen ist aber eine ganz verschiedene, jenachdem die Determinante \(D\) eine positive (nicht quadratische) oder negative Zahl ist. Während im letztern Falle die zweite Untersuchung sehr einfach ist, bereitet sie im erstern grössere Schwierigkeit. In ihm zerfallen die reduzierten Formen in eine endliche Anzahl von Perioden benachbarter Formen, welche also einander äquivalent sind; dass aber Formen verschiedener Perioden es nicht sind, erfordert einen umständlicheren Beweis, der, wenn man auf sein Wesen sieht, auf die Entwicklung der sogenannten »Wurzeln« einer reduzierten Form in einen Kettenbruch begründet ist, wie später \textsc{Dirichlet}\begingroup\renewcommand{\thefootnote}{**)}\footnote{Vergl. \textit{Vorlesungen über Zahlentheorie}, herausg. von \textsc{Dedekind}, 4. Aufl. 1894 § 79—82.}\endgroup\ sehr schön und einfach klargelegt hat. Wenn nun auf dem angedeuteten Wege eine Transformation der Form \(f\) in eine der Formen \(F\) ermittelt ist, so lassen sie sich \textls[100]{sämtlich} aufstellen, sobald man alle Auflösungen der von \textsc{Euler}\begingroup\renewcommand{\thefootnote}{***)}\footnote{Comment. Acad. Petrop. \textbf{6} (1732/3) 1738, S. 175, Opera omnia, Ser. I, vol. 2, S. 6; Novi comment. Acad. Petrop. \textbf{11} (1765) 1767, S. 28; \textit{Vollständige Anleitung zur Algebra} (1770), II. Teil, II. Abschnitt, Kap. 7; Opuscula analytica I (1783), S. 310.}\endgroup\ als \textsc{Pell}schen bezeichneten Gleichung
\[
t^2-Du^2=1
\]
anzugeben weiss. Diesen Teil der Aufgabe, der für eine negative Determinante sich unmittelbar erledigt, löste für den Fall einer positiven Determinante schon \textsc{Lagrange}\begingroup\renewcommand{\thefootnote}{\ensuremath{\dagger})}\footnote{Miscell. Taur. \textbf{4} (1766/9), S. 19; Oeuvres I, S. 669; Mém. de l’Acad. Berlin (1767) 1769, \textbf{23}, S. 165, Oeuvres II, S. 375; \textit{Additions aux Elémens d’Algèbre par L. Euler}, § II, VIII, Oeuvres VII, S. 45, 157.}\endgroup, der auch zuerst die Existenz einer Auflösung nachwies, auf Grund der Kettenbruchentwicklung für die Quadratwurzel aus \(D\). Die für den
\newpage\SourceLeaf{28}\ArithmeticHead{21}{R}{BINÄRE QUADRATISCHE FORMEN.}
\noindent gleichen Zweck von \textsc{Gauss} gegebene Analyse, von der er sagt, sie sei ex principiis omnino diversis petita, und welche durch Verbindung der Transformation einer reduzierten Form in die erste in ihrer Periode ihr gleiche mit der identischen Transformation zum Ziele gelangt, kommt gleichwohl, da die Äquivalenz der aufeinanderfolgenden reduzierten Formen mit der Äquivalenz ihrer Wurzeln identisch, die letztere aber mit jener Kettenbruchentwicklung aufs engste verbunden ist, im Wesen mit der von \textsc{Lagrange} überein.

So haben wir die Ergebnisse der von \textsc{Gauss} in streng logischer Verkettung und mit grösster Vollständigkeit unter Berücksichtigung aller Nebenumstände, wie der Formen mit quadratischer oder der Null gleicher Determinante usw., entwickelten Elemente der Theorie in der Hauptsache geschildert. Als deren Abschluss können wir die Verteilung aller Formen mit einer gegebenen Determinante auf eine endliche Anzahl von Klassen unter einander äquivalenter Formen, die übrigens schon \textsc{Lagrange}\footnote{Nouv. Mém. de l’Acad. Berlin (1773) 1775, S. 265; Oeuvres III, S. 695.}\ zugehört, betrachten. Aus der allgemeinen Theorie aber fliessen nun als aus ihrer eigentlich wissenschaftlichen Quelle die schon von \textsc{Fermat}\begingroup\renewcommand{\thefootnote}{**)}\footnote{Oeuvres II, S. 403 ff.}\endgroup\ erkannten Sätze von der Darstellbarkeit einer Primzahl von der Form \(4n+1\), von einer der Formen \(8n+1\) oder \(8n+3\), und von der Form \(6n+1\) bezw. in den Formen \(x^2+y^2,\ x^2+2y^2,\ x^2+3y^2\), von denen der erste durch \textsc{Euler}\begingroup\renewcommand{\thefootnote}{***)}\footnote{Novi Comment. Acad. Petrop. \textbf{1} (1747/8) 1750, S. 20; \textbf{4} (1752/3) 1758, S. 3; \textbf{5} (1754/5) 1760, S. 3, Opera Omnia, Ser. I, vol. 2, S. 62; 295; 328.}\endgroup, der zweite durch \textsc{Lagrange}\begingroup\renewcommand{\thefootnote}{\ensuremath{\dagger})}\footnote{Nouv. Mém. de l’Acad. Berlin (1775) 1777, S. 323, Oeuvres III, S. 759. Vgl. auch \textsc{Euler}, Novi Comment. Acad. Petrop. \textbf{6} (1756/7) 1761, S. 185, Opera Omnia, Ser. I, vol. 2, S. 459.}\endgroup, der dritte wieder durch \textsc{Euler}\begingroup\renewcommand{\thefootnote}{\ensuremath{\dagger\dagger})}\footnote{Novi Comment. Acad. Petrop. \textbf{8} (1760/1) 1763, S. 105, Opera Omnia, Ser. I, vol. 2, S. 556.}\endgroup\ zuerst bewiesen worden war. Den ersten derselben hat \textsc{Gauss} an einer späteren Stelle der D. A. (art. 265) noch einmal aus tieferen Prinzipien wieder hergeleitet.
